% ============================================================
%  Supplementary Material
%  Quantifying LLM Non-Determinism in Evidence Synthesis
%  Target: Research Synthesis Methods
% ============================================================
\documentclass[11pt]{article}

\usepackage[T1]{fontenc}
\usepackage{newtxtext}
\usepackage[cmintegrals]{newtxmath}
\usepackage{newtxtt}
\usepackage{graphicx}
\usepackage{booktabs}
\usepackage{geometry}
\usepackage{listings}
\usepackage{xcolor}
\usepackage{hyperref}
\usepackage{longtable}

\geometry{margin=2.5cm}

\lstset{
  basicstyle=\ttfamily\small,
  breaklines=true,
  frame=single,
  backgroundcolor=\color{gray!5},
  keywordstyle=\color{blue!70},
  stringstyle=\color{green!50!black},
  showstringspaces=false,
  tabsize=2,
}

\title{Supplementary Material\\[0.5em]
\large When the Same Question Gets Different Answers:\\
Quantifying LLM Non-Determinism in Evidence Synthesis}
\author{Lucas Rover \and Yara de Souza Tadano}
\date{}

\begin{document}

\maketitle
\tableofcontents
\newpage

% ============================================================
\section{PubMed Search Strategy}
\label{sec:search}

All abstracts were retrieved from PubMed via the NCBI E-utilities API (\texttt{esearch} + \texttt{efetch}) on February 11, 2026.
Two complementary queries were used:

\subsection{Broad Query (573 records)}

\begin{lstlisting}[language={},caption={PubMed broad search query}]
("Particulate Matter"[MeSH] OR "PM2.5" OR "fine particulate matter"
 OR "fine particles" OR "ambient particulate matter")
AND ("Hospitalization"[MeSH] OR "hospital admission"
 OR "hospital admissions" OR "emergency department"
 OR "ED visit" OR "ED visits")
AND ("Respiratory Tract Diseases"[MeSH] OR "respiratory"
 OR "asthma" OR "COPD" OR "pneumonia" OR "bronchitis"
 OR "bronchiolitis")
AND ("time series" OR "time-series" OR "case-crossover"
 OR "distributed lag" OR "Poisson regression" OR "GAM"
 OR "generalized additive model")
\end{lstlisting}

\subsection{Design-Filtered Query (222 records)}

\begin{lstlisting}[language={},caption={PubMed design-filtered query}]
("PM2.5" OR "fine particulate matter"
 OR "particulate matter 2.5")
AND ("respiratory" OR "asthma" OR "COPD" OR "pneumonia")
AND ("hospitalization" OR "hospital admission"
 OR "emergency department")
AND ("time series" OR "time-series" OR "case-crossover")
\end{lstlisting}

\subsection{Exclusion Candidates (475 records)}

A third query targeted articles clearly outside scope (animal studies, in-vitro, occupational exposure) to populate the ``clearly exclude'' category.

\subsection{Corpus Assembly}

From the merged candidate pool (1,021 unique records after deduplication):
\begin{itemize}
  \item 100 abstracts meeting all six inclusion criteria were labeled ``clearly include''
  \item 100 abstracts with unambiguous exclusion reasons were labeled ``clearly exclude''
  \item 300 abstracts with borderline characteristics (e.g., mixed pollutants, ambiguous outcomes) were labeled ``ambiguous''
\end{itemize}

The 200 ``clear'' abstracts were independently labeled by two reviewers with discordance resolution.
The 300 ``ambiguous'' abstracts were classified using study-design indicators and keyword matching (see Labeling Guide, Section~\ref{sec:labeling}).

% ============================================================
\section{Prompt Templates}
\label{sec:prompts}

The exact prompts used for both pipeline stages are reproduced below.
These prompts were held constant across all 6 models and all 120 experimental runs.

\subsection{Stage A: Abstract Screening Prompt}

\begin{lstlisting}[language={},caption={Screening prompt template (\texttt{configs/prompts/screening.txt})}]
You are an expert epidemiologist conducting a systematic review on
the association between PM2.5 exposure and respiratory
hospitalizations.

Given the following title and abstract of a scientific publication,
determine whether it should be INCLUDED or EXCLUDED from the review
based on the criteria below.

## Inclusion Criteria
1. Original epidemiological study (not a review, editorial, or
   commentary)
2. Exposure: PM2.5 or fine particulate matter (<=2.5 um)
3. Outcome: respiratory hospitalization or emergency department visit
4. Study design: time-series or case-crossover
5. Reports a relative risk (RR), odds ratio (OR), or equivalent
   effect measure with confidence interval
6. Published in English

## Exclusion Criteria
- Reviews, meta-analyses, editorials, or commentaries
- Exposure other than PM2.5 (unless PM2.5 is one of multiple
  exposures analyzed)
- Outcome other than respiratory hospitalization/ED visit
- No quantitative effect estimate reported
- Animal or in-vitro studies

## Input
Title: {title}
Abstract: {abstract}

## Instructions
Respond with a JSON object only. Do not include any text outside
the JSON.

{
  "decision": "include" | "exclude" | "uncertain",
  "confidence": 0.0-1.0,
  "rationale": "brief explanation of decision",
  "exposure": "PM2.5" | "other" | "not_specified",
  "outcome": "respiratory_hospitalization" | "respiratory_ED"
             | "other" | "not_specified",
  "study_design": "time_series" | "case_crossover" | "cohort"
                  | "other" | "not_specified",
  "has_effect_estimate": true | false
}
\end{lstlisting}

\subsection{Stage B: Data Extraction Prompt}

\begin{lstlisting}[language={},caption={Extraction prompt template (\texttt{configs/prompts/extraction.txt})}]
You are an expert epidemiologist extracting quantitative data from a
scientific publication on PM2.5 and respiratory hospitalizations for
a systematic review and meta-analysis.

Given the following title and abstract, extract the structured data
below. If a field is not reported in the abstract, use null.

## Input
Title: {title}
Abstract: {abstract}

## Instructions
Extract ALL reported effect estimates. If multiple estimates are
reported (e.g., different lags, subgroups), extract each as a
separate entry in the "estimates" array.

Respond with a JSON object only. Do not include any text outside
the JSON.

{
  "study_id": "first_author_year",
  "study_location": "city, country",
  "study_period": "YYYY-YYYY",
  "study_design": "time_series" | "case_crossover" | "other",
  "population": "general" | "elderly" | "children" | "other",
  "sample_size": null | integer,
  "estimates": [
    {
      "effect_measure": "RR" | "OR" | "HR" | "IRR" | "other",
      "effect_estimate": float,
      "ci_lower": float,
      "ci_upper": float,
      "ci_level": 95,
      "exposure_increment": "per 10 ug/m3" | "per IQR" | "other",
      "lag": "lag0" | "lag1" | "lag0-1" | "other",
      "outcome_specific": "all_respiratory" | "pneumonia"
                         | "COPD" | "asthma" | "other",
      "covariates": ["temperature", "humidity", "trend",
                     "seasonality", "..."]
    }
  ]
}
\end{lstlisting}

% ============================================================
\section{JSON Schemas for Output Validation}
\label{sec:schemas}

All LLM outputs were validated against JSON Schema (Draft-07) definitions.
Outputs failing validation were flagged but retained to ensure complete provenance.

\subsection{Screening Output Schema}

Seven required fields, no additional properties permitted:
\texttt{decision} (enum: include/exclude/uncertain),
\texttt{confidence} (0.0--1.0),
\texttt{rationale} (string, min 10 chars),
\texttt{exposure}, \texttt{outcome}, \texttt{study\_design} (categorical enums),
and \texttt{has\_effect\_estimate} (boolean).

\subsection{Extraction Output Schema}

Required fields: \texttt{study\_id} (string), \texttt{estimates} (array).
Optional metadata: \texttt{study\_location}, \texttt{study\_period}, \texttt{study\_design}, \texttt{population}, \texttt{sample\_size}.
Each estimate requires: \texttt{effect\_measure}, \texttt{effect\_estimate}, \texttt{ci\_lower}, \texttt{ci\_upper} (all with range constraints 0.1--10.0).

Full schemas are available in the repository at \texttt{configs/schemas/}.

% ============================================================
\section{Gold Standard Labeling Protocol}
\label{sec:labeling}

\subsection{Screening Labels}

The 200 ``clear'' abstracts (100 include, 100 exclude) were independently labeled by two reviewers following a structured protocol.
Discordance was resolved by a third reviewer.
Inter-rater agreement targets: Cohen's $\kappa \geq 0.80$, discordance rate $< 15\%$.

The 300 ``ambiguous'' abstracts were classified using keyword matching and study-design indicators.
These labels serve as reference values for accuracy computation but do not affect the primary metric (EMR), which measures self-consistency regardless of ground truth.

\subsection{Extraction Labels}

For the 100 ``clearly include'' articles, extraction templates were created following standardized rules:
\begin{itemize}
  \item Prefer the primary/main estimate if multiple are reported
  \item Prefer all-respiratory over specific conditions (asthma, COPD)
  \item Prefer single-pollutant model over multi-pollutant
  \item Prefer lag 0--1 if multiple lags reported without a ``best'' designation
  \item Record per 10~$\mu$g/m$^3$ -- convert if different increment used
  \item If abstract reports \% change, convert: RR $= 1 + (\%\text{change} / 100)$
\end{itemize}

Discordance threshold: $>0.005$ for effect estimates, $>0.01$ for CI bounds.

% ============================================================
\section{Example Outputs: Exact Match vs.\ Near-Miss}
\label{sec:examples}

To illustrate the nature of extraction non-determinism, we present two representative cases.

\subsection{Example 1: Exact Match (Local Model -- Gemma 2 9B)}

All 10 runs produced identical output (EMR = 1.000):

\begin{lstlisting}[language={},caption={Gemma 2 9B extraction output (identical across all 10 runs)}]
{
  "study_id": "chen_2017",
  "study_location": "Shanghai, China",
  "study_period": "2013-2015",
  "study_design": "time_series",
  "population": "general",
  "sample_size": null,
  "estimates": [
    {
      "effect_measure": "RR",
      "effect_estimate": 1.023,
      "ci_lower": 1.005,
      "ci_upper": 1.042,
      "ci_level": 95,
      "exposure_increment": "per 10 ug/m3",
      "lag": "lag0-1",
      "outcome_specific": "all_respiratory",
      "covariates": ["temperature", "humidity", "day_of_week",
                     "long_term_trend", "seasonality"]
    }
  ]
}
\end{lstlisting}

\subsection{Example 2: Near-Miss (Cloud API -- Claude Sonnet 4.5)}

Runs 1--8 and 10 produced one output; Run 9 differed in the \texttt{estimates} array:

\begin{lstlisting}[language={},caption={Claude Sonnet 4.5 -- Run 9 extracted 2 estimates instead of 1}]
{
  "study_id": "chen_2017",
  "study_location": "Shanghai, China",
  "study_period": "2013-2015",
  "study_design": "time_series",
  "population": "general",
  "sample_size": null,
  "estimates": [
    {
      "effect_measure": "RR",
      "effect_estimate": 1.023,
      "ci_lower": 1.005,
      "ci_upper": 1.042,
      "ci_level": 95,
      "exposure_increment": "per 10 ug/m3",
      "lag": "lag0-1",
      "outcome_specific": "all_respiratory",
      "covariates": ["temperature", "humidity", "day_of_week",
                     "long_term_trend", "seasonality"]
    },
    {
      "effect_measure": "RR",
      "effect_estimate": 1.031,
      "ci_lower": 1.008,
      "ci_upper": 1.054,
      "ci_level": 95,
      "exposure_increment": "per 10 ug/m3",
      "lag": "lag2",
      "outcome_specific": "pneumonia",
      "covariates": ["temperature", "humidity", "day_of_week"]
    }
  ]
}
\end{lstlisting}

\noindent \textbf{Key observation:} The metadata fields (\texttt{study\_id}, \texttt{study\_location}, \texttt{study\_design}) are identical across all 10 runs (consistent with BERTScore $F_1 = 0.997$).
The variation lies in the \texttt{estimates} array---Run 9 extracted an additional sub-group estimate (pneumonia at lag 2) that the other 9 runs did not.
This single additional estimate causes the hash-based EMR to register a non-match, while the BERTScore over metadata text remains near-perfect.
This exemplifies the ``semantic paradox'' reported in Section 4.5 of the main paper.

% ============================================================
\section{Model Configurations}
\label{sec:configs}

\begin{table}[htbp]
\centering
\caption{Complete model configurations for all six LLMs.}
\small
\begin{tabular}{llccccl}
\toprule
\textbf{Model} & \textbf{Provider} & \textbf{Temp.} & \textbf{Seed} & \textbf{Max Tokens} & \textbf{Timeout} & \textbf{Version} \\
\midrule
LLaMA 3 8B     & Ollama   & 0.0 & 42   & 4096  & 600s & \texttt{llama3:8b} \\
Mistral 7B     & Ollama   & 0.0 & 42   & 4096  & 600s & \texttt{mistral:7b} \\
Gemma 2 9B     & Ollama   & 0.0 & 42   & 4096  & 600s & \texttt{gemma2:9b} \\
\midrule
Claude Sonnet 4.5 & Anthropic & 0.0 & --- & 4096 & 90s & \texttt{claude-sonnet-4-5-20250929} \\
Gemini 2.5 Pro & Google   & 0.0 & 42   & 8192$^\dagger$ & 90s & \texttt{gemini-2.5-pro} \\
GPT-4.1        & OpenAI   & 0.0 & 42   & 4096  & 90s & \texttt{gpt-4.1} \\
\bottomrule
\multicolumn{7}{l}{\footnotesize $^\dagger$Gemini requires higher \texttt{maxOutputTokens} because thinking tokens consume the output budget.}
\end{tabular}
\end{table}

% ============================================================
\section{Provenance Protocol Details}
\label{sec:provenance}

\subsection{Call-Level Hashing}

Each LLM call produces a \textbf{call hash} computed as:

\begin{lstlisting}[language=Python,caption={Call hash computation (\texttt{src/provenance/hasher.py})}]
import hashlib, json

def compute_call_hash(prompt, input_text, model_id, temperature,
                      seed=None):
    fields = {
        "prompt": prompt,
        "input_text": input_text,
        "model_id": model_id,
        "temperature": temperature,
        "seed": seed,
    }
    canonical = json.dumps(fields, sort_keys=True,
                           ensure_ascii=True)
    return hashlib.sha256(canonical.encode()).hexdigest()
\end{lstlisting}

The use of \texttt{sort\_keys=True} and \texttt{ensure\_ascii=True} ensures deterministic serialization across platforms and Python versions.

\subsection{Output Hashing}

The raw LLM response text is hashed directly:

\begin{lstlisting}[language=Python]
def compute_output_hash(output_text):
    return hashlib.sha256(output_text.encode()).hexdigest()
\end{lstlisting}

\subsection{Run Cards}

Each run produces a structured JSON run card containing:
\begin{itemize}
  \item Run metadata: model ID, provider, stage, run number
  \item Configuration: temperature, seed, max tokens
  \item Execution: timestamps, call counts (total/success/fail), duration
  \item Provenance: aggregate output hash (SHA-256 of sorted individual hashes joined by pipe delimiter)
\end{itemize}

\subsection{Verification Protocol}

If two runs of the same model produce different aggregate hashes, the per-call output hashes are compared to identify exactly which abstracts changed.
This enables targeted analysis of non-determinism at the individual-abstract level.

% ============================================================
\section{Corpus Statistics}
\label{sec:corpusstats}

\begin{table}[htbp]
\centering
\caption{Corpus descriptive statistics.}
\begin{tabular}{lr}
\toprule
\textbf{Statistic} & \textbf{Value} \\
\midrule
Total abstracts & 500 \\
Unique journals & 148 \\
Publication range & 1994--2026 \\
Mean abstract length (characters) & 1,873 \\
Median abstract length (characters) & 1,812 \\
Min abstract length (characters) & 423 \\
Max abstract length (characters) & 4,291 \\
\midrule
Clear include & 100 (20\%) \\
Clear exclude & 100 (20\%) \\
Ambiguous & 300 (60\%) \\
\bottomrule
\end{tabular}
\end{table}

% ============================================================
\section{Experimental Timeline}
\label{sec:timeline}

All 120 experimental runs were conducted within a 3-week window to minimize the risk of silent API model updates:

\begin{table}[htbp]
\centering
\caption{Experiment execution timeline.}
\begin{tabular}{llll}
\toprule
\textbf{Model} & \textbf{Screening} & \textbf{Extraction} & \textbf{Duration} \\
\midrule
LLaMA 3 8B      & Feb 11--12 & Feb 12--13 & $\sim$91h \\
Mistral 7B       & Feb 13--15 & Feb 15--16 & $\sim$77h \\
Gemma 2 9B       & Feb 16--17 & Feb 17     & $\sim$46h \\
Claude Sonnet 4.5 & Feb 18    & Feb 18--19 & $\sim$12h \\
Gemini 2.5 Pro   & Feb 20    & Feb 20--21 & $\sim$8h \\
GPT-4.1          & Feb 22    & Feb 22--Mar 2 & $\sim$6h \\
\bottomrule
\end{tabular}
\end{table}

% ============================================================
\section{Software and Dependencies}
\label{sec:software}

\begin{itemize}
  \item \textbf{Python}: 3.14.3
  \item \textbf{OS}: macOS Darwin 24.6.0 (Apple M4, 24 GB RAM)
  \item \textbf{Ollama}: v0.15.5 (local model serving)
  \item \textbf{Key dependencies}: numpy $\geq$ 1.26, pandas $\geq$ 2.1, scipy $\geq$ 1.12, jsonschema $\geq$ 4.20, bert-score $\geq$ 0.3.13, transformers $\geq$ 4.36
  \item \textbf{No external SDKs}: All API calls use Python's \texttt{urllib.request} to ensure uniform request handling
  \item \textbf{Test suite}: 45 tests (pytest $\geq$ 8.0), 100\% pass rate
  \item \textbf{License}: MIT
  \item \textbf{Repository}: \url{https://github.com/Roverlucas/llm-evidence-synthesis-reproducibility}
\end{itemize}

% ============================================================
\section{Complete Bootstrap Confidence Intervals}
\label{sec:bootstrap}

All confidence intervals were computed using the percentile bootstrap method with 10,000 resamples.
For models with EMR = 1.000, the bootstrap trivially returns [1.000, 1.000]; the rule-of-three upper bound on the non-match rate is $3/n$ (0.6\% for $n = 500$ screening, 3.0\% for $n = 100$ extraction).

\begin{table}[htbp]
\centering
\caption{Complete EMR with 95\% bootstrap confidence intervals for all model-stage combinations.}
\small
\begin{tabular}{llccc}
\toprule
\textbf{Model} & \textbf{Stage} & \textbf{EMR} & \textbf{95\% CI} & \textbf{$n$} \\
\midrule
LLaMA 3 8B       & Screening  & 1.000 & [1.000, 1.000]$^*$ & 480 \\
LLaMA 3 8B       & Extraction & 1.000 & [1.000, 1.000]$^*$ & 96 \\
Mistral 7B        & Screening  & 1.000 & [1.000, 1.000]$^*$ & 500 \\
Mistral 7B        & Extraction & 1.000 & [1.000, 1.000]$^*$ & 100 \\
Gemma 2 9B        & Screening  & 1.000 & [1.000, 1.000]$^*$ & 500 \\
Gemma 2 9B        & Extraction & 1.000 & [1.000, 1.000]$^*$ & 100 \\
\midrule
Claude Sonnet 4.5 & Screening  & 0.974 & [0.960, 0.986] & 500 \\
Claude Sonnet 4.5 & Extraction & 0.050 & [0.010, 0.090] & 100 \\
Gemini 2.5 Pro    & Screening  & 0.936 & [0.914, 0.956] & 500 \\
Gemini 2.5 Pro    & Extraction & 0.200 & [0.120, 0.280] & 100 \\
GPT-4.1           & Screening  & 0.970 & [0.954, 0.984] & 500 \\
GPT-4.1           & Extraction & 0.150 & [0.080, 0.220] & 100 \\
\bottomrule
\multicolumn{5}{l}{\footnotesize $^*$Rule-of-three upper bound on non-match rate: $3/n$ (screening: 0.6\%, extraction: 3.0\%).}
\end{tabular}
\end{table}

\end{document}
